\newpage

% ######################################################################
% PART VIII: THEORY DEEP DIVE
% ######################################################################
\part{Theory Deep Dive --- Interview Gold}

% ######################################################################
% SECTION 37: TRANSFORMER ARCHITECTURE
% ######################################################################
\section{Transformer Architecture --- LLM Ki Jaan}

\begin{storybox}
Interviewer puchta hai: ``LLM ke andar kya hai?'' --- is section mein transformer ka pura breakdown hai: attention, embeddings, layers, training. \textbf{Bina transformer samjhe LLM interview answer adhura hai.} Is project ke models (GPT-OSS-120B) transformer-based hain --- isliye ye theory directly relevant hai.
\end{storybox}

\subsection{Transformer Kya Hai --- 2017 Ka Revolution}

\begin{conceptbox}
Transformer: \textbf{``Attention Is All You Need''} paper (Vaswani et al., 2017) se aaya architecture. Pehle RNNs/LSTMs text process karte the sequentially (word-by-word) --- slow, aur long-range dependencies bhool jate the. Transformer ka breakthrough: \textbf{saare words ko parallel process karo}, aur \textbf{attention} se words ke beech relationships capture karo.

\textbf{2 core ideas:}
\begin{enumerate}[label=\arabic*.]
\item \textbf{Self-attention} --- har word ko baaki sab words se context milta hai.
\item \textbf{Parallelization} --- training bahut fast (isliye scale possible hua).
\end{enumerate}
\end{conceptbox}

\subsection{Embedding Layer --- Words Se Numbers}

\begin{mathbox}
Har token ko ek vector mile: $w_i \in \mathbb{R}^{d}$ jahan $d$ = embedding dimension (GPT-OSS class models mein 4096+).

\textbf{Positional encoding:} Transformer ko word ka order pata hona chahiye (``dog bites man'' vs ``man bites dog''). Position $pos$ ke liye:
\[
PE_{(pos, 2i)} = \sin\left(\frac{pos}{10000^{2i/d}}\right), \quad
PE_{(pos, 2i+1)} = \cos\left(\frac{pos}{10000^{2i/d}}\right)
\]
Ye sine/cosine pattern model ko relative positions samajhne deta hai.
\end{mathbox}

\subsection{Self-Attention --- The Heart}

\begin{mathbox}
\textbf{Scaled Dot-Product Attention:}
\[
\text{Attention}(Q, K, V) = \text{softmax}\left(\frac{QK^T}{\sqrt{d_k}}\right)V
\]

\textbf{Step-by-step:}
\begin{enumerate}[label=\arabic*.]
\item Har token ke liye 3 vectors: \textbf{Query} (main poochta hoon), \textbf{Key} (tum kya ho), \textbf{Value} (tum kya dete ho).
\item $Q \cdot K^T$: har token ka har token se \textbf{similarity score}.
\item $\sqrt{d_k}$ se divide: scores stable (large dims mein softmax saturation se bachao).
\item Softmax: scores $\rightarrow$ probabilities (kitna attention kis par).
\item Probabilities $\times$ Values: weighted sum = token ka context-aware representation.
\end{enumerate}

\textbf{Intuition:} Sentence ``The cat sat on the mat because it was tired'' mein \texttt{it} ko \texttt{cat} se connect karne ke liye attention hi mechanism hai.
\end{mathbox}

\subsection{Multi-Head Attention}

\begin{conceptbox}
Ek attention head ek \textbf{relationship type} capture karta hai (syntax, co-reference, sentiment). Multi-head: \textbf{8-64 heads parallel} --- har head alag relationship family. Outputs concatenate hote hain.

\[
\text{MultiHead}(Q,K,V) = \text{Concat}(head_1, ..., head_h)W^O
\]
Is project ke GPT-OSS-120B mein ye multi-head mechanism hi har agent prompt ka context samajhta hai.
\end{conceptbox}

\subsection{Transformer Block Stack}

\begin{conceptbox}
\textbf{Har block (layer) mein:}
\begin{enumerate}[label=\arabic*.]
\item Multi-head self-attention
\item Residual connection (skip) + LayerNorm
\item Feed-forward network (2 linear layers + ReLU/GELU)
\item Residual connection + LayerNorm
\end{enumerate}

GPT-OSS-120B mein \textbf{dozens of such blocks stacked} hain. Depth = capability: deeper model complex patterns seekhta hai. Output layer: next-token probability distribution over vocabulary (\textasciitilde100K tokens).
\end{conceptbox}

\begin{interviewbox}
\textbf{Q: ``GPT ka 'Generative Pre-trained Transformer' kya matlab hai?''}

\textbf{Answer:} (1) \textbf{Generative} --- next token predict karta hai; (2) \textbf{Pre-trained} --- pehle massive text par self-supervised training; (3) \textbf{Transformer} --- architecture. Training 2-stage: pre-training (next-word prediction, trillions tokens) $\rightarrow$ fine-tuning (instruction following, RLHF). Is project ka model inference par chalata hai --- trained model se sirf generate karte hain.
\end{interviewbox}

\subsection{Training vs Inference}

\begin{longtable}{p{4cm} p{5.2cm} p{5.2cm}}
\toprule
\textbf{Aspect} & \textbf{Training} & \textbf{Inference} \\
\midrule
Goal & Learn weights & Generate output \\
Data & Trillions tokens & User prompt \\
Compute & GPUs, weeks-months & Single forward pass \\
Cost & Millions USD & Cents per request \\
Updates & Backpropagation & No updates \\
Is project mein? & Provider side (Groq) & \textbf{Yes --- every API call} \\
\bottomrule
\end{longtable}

\begin{memorybox}
\textbf{Interview pearl:} ``Aapka project LLM \textbf{train} nahi karta --- trained model ko \textbf{use} karta hai (inference via API). Model chunna, prompt design karna, aur evaluation karna --- ye humara kaam hai. Fine-tuning tab chahiye jab domain-specific style/format chahiye jo prompting se nahi aata.''
\end{memorybox}

% ######################################################################
% SECTION 38: ADVANCED RAG
% ######################################################################
\section{Advanced RAG --- Naive Se Production Tak}

\subsection{Naive RAG ki Problems}

\begin{conceptbox}
Naive RAG (retrieve $\rightarrow$ stuff $\rightarrow$ generate) ki 4 problems:
\begin{enumerate}[label=\arabic*.]
\item \textbf{Retrieval quality} --- top-k chunks irrelevant ho sakte hain (query phrasing par sensitive).
\item \textbf{Lost in the middle} --- LLM context ke middle ko ignore karta hai (position bias).
\item \textbf{Chunk boundaries} --- context beech mein kat sakta hai.
\item \textbf{No evaluation} --- retrieval/answer quality measured nahi.
\end{enumerate}
\end{conceptbox}

\subsection{Advanced RAG Techniques}

\begin{longtable}{p{4.2cm} p{9.8cm}}
\toprule
\textbf{Technique} & \textbf{Kya karta hai} \\
\midrule
Query rewriting & LLM se query ko expand/compress karna (is project: researcher ke 3 sub-queries) \\
Hybrid search & Vector + BM25 keyword search combine (dense + sparse) \\
Re-ranking & Cross-encoder se top-50 $\rightarrow$ top-5 refine \\
Context compression & LLM se retrieved chunks ko compress (LLMLingua) \\
Parent-child chunks & Chhote chunks retrieve, bade parent chunks LLM ko do \\
Metadata filtering & Date/source/type filters pehle se \\
RAPTOR & Hierarchical summaries build karke multi-level retrieve \\
Self-RAG & Model khud decide kare: retrieve karna hai ya nahi \\
\bottomrule
\end{longtable}

\begin{interviewbox}
\textbf{Q: ``Is project mein kaunsa advanced technique use hua?''}

\textbf{Answer:} \textbf{Query expansion} (researcher 3 diverse sub-queries banata hai --- naive single-query ka improved version) aur \textbf{multi-source retrieval} (Tavily web + research\_docs + past\_research --- three diverse corpora). Production version mein main re-ranking (cross-encoder) aur hybrid search add karta. Ye progression batana strong answer deta hai.
\end{interviewbox}

\subsection{RAG Evaluation}

\begin{conceptbox}
\textbf{2-layer evaluation:}
\begin{enumerate}[label=\arabic*.]
\item \textbf{Retrieval quality:} Recall@k, MRR --- ``kya relevant chunks top-k mein aaye?''
\item \textbf{Generation quality:} Faithfulness (answer sources se hai?), Relevance, Hallucination rate.
\end{enumerate}
Is project ka benchmark \textbf{verifiability} dimension exactly generation-layer check karta hai --- citations + facts-vs-opinions. Interview mein ye connection batao: ``benchmark ki verifiability = RAG faithfulness ka proxy.''
\end{conceptbox}

% ######################################################################
% SECTION 39: DEBUGGING SESSION NARRATIVE
% ######################################################################
\section{Debugging Session --- Ek Real Kahani}

\begin{storybox}
Ye section ek \textbf{real debugging session} ki kahani hai jo maine is project par ki --- sab kuch sach hai jo is document mein likha hai. Interview mein story-telling format mein answer dena (STAR) powerful hota hai --- ye ek ready-made example hai.
\end{storybox}

\subsection{Bug 1: MCP Server Crash on Startup}

\begin{interviewbox}
\textbf{Symptom:} \texttt{mcp\_server.py} import karte hi crash: \texttt{FastMCP() got an unexpected keyword argument 'description'}.

\textbf{Diagnosis:} fastmcp library version 3.4.6 mein \texttt{description} parameter rename hua \texttt{instructions}. Code purane version ke liye likha tha --- dependency update ke baad toot gaya.

\textbf{Fix:} \texttt{FastMCP(description=...)} $\rightarrow$ \texttt{FastMCP(instructions=...)}. Verified: sab 5 tools + resource register hote hain.

\textbf{Lesson:} Pin dependencies (requirements.txt mein versions) + upgrade ke baad smoke test. Ye ``dependency drift'' problem har production system mein aati hai.
\end{interviewbox}

\subsection{Bug 2: Model 404 on Groq}

\begin{interviewbox}
\textbf{Symptom:} Pipeline mein har LLM call 404 error --- \texttt{llama-3.3-70b-versatile} not found on account.

\textbf{Diagnosis:} Groq account par wo model available hi nahi tha. Mainne \texttt{GET /v1/models} call karke asli available models dekhe --- sirf prompt-guard classifiers + openai/gpt-oss + qwen models mile.

\textbf{Fix:} Live tests kiye teen models par --- gpt-oss-120b ne clean output diya (qwen reasoning model \texttt{<think>} blocks deta tha jo report prose kharab karta). \texttt{MODEL\_NAME=openai/gpt-oss-120b} set kiya config + .env + README + notes mein.

\textbf{Lesson:} ``Documented model'' ka matlab available model nahi --- hamesha live API se verify karo. Agents ki graceful fallback ne system ko crash hone se bachaya.
\end{interviewbox}

\subsection{Bug 3: History Ordering Nondeterministic}

\begin{interviewbox}
\textbf{Symptom:} Test fail --- do benchmark rows ek hi timestamp ke saath, ordering flaky.

\textbf{Diagnosis:} Django test env \textbf{time freezes} --- dono rows ka \texttt{created\_at} identical, \texttt{order\_by("-created\_at")} tie par nondeterministic.

\textbf{Fix:} \texttt{order\_by("-created\_at", "-id")} --- deterministic tiebreak.

\textbf{Lesson:} Production mein bhi same-timestamp rows possible hain (bulk inserts) --- \textbf{deterministic ordering} hamesha secondary key rakho.
\end{interviewbox}

\subsection{Bug 4: venv Silent Fallback}

\begin{interviewbox}
\textbf{Symptom:} Django \texttt{No module named 'django'} jabki manage.py pehle chalta tha.

\textbf{Diagnosis:} Git Bash mein \texttt{source .venv/bin/activate} (Linux-style path) Windows venv par \textbf{silently fail} karta hai --- Python global se chalta tha. Naye environment mein django missing tha.

\textbf{Fix:} \texttt{.venv/Scripts/python.exe} use karo + full requirements.txt install.

\textbf{Lesson:} ``Aur environment setup documentation likho'' --- ye exact issue run.md mein documented hai. Cross-platform tooling ki path differences real problems deti hain.
\end{interviewbox}

\subsection{Bug 5: LaTeX Caption Underscores}

\begin{interviewbox}
\textbf{Symptom:} Ye notes document compile karte waqt \texttt{Missing \$ inserted} errors --- sirf kuch listings par.

\textbf{Diagnosis:} Captions mein \texttt{rag/vector\_store.py} jaise filenames --- text mode mein \texttt{\_} = math subscript. Isliye sirf underscore wale captions fail hote the (vector\_store, fact\_checker, mcp\_server, 1\_Benchmark) --- baaki (schema.py, workflow.py) clean the.

\textbf{Fix:} Captions mein \texttt{\_} escape kiya.

\textbf{Lesson:} \textbf{Reproduce minimally} --- maine failing listing ko alag file mein isolate kiya aur byte-level diff kiya. Debugging ka systematic approach: hypothesis $\rightarrow$ minimal repro $\rightarrow$ root cause $\rightarrow$ fix.
\end{interviewbox}
